\documentclass[11pt]{amsart}

\usepackage{amssymb}
\usepackage[hidelinks]{hyperref}

\newtheorem{theorem}{Theorem}[section]
\newtheorem{lemma}[theorem]{Lemma}
\newtheorem{proposition}[theorem]{Proposition}
\newtheorem{corollary}[theorem]{Corollary}
\theoremstyle{definition}
\newtheorem{definition}[theorem]{Definition}
\theoremstyle{remark}
\newtheorem{remark}[theorem]{Remark}

\title{{{TITLE}}}

\author{First Author}
\address{Department of Mathematics, University One, City, Country}
\email{first@example.com}

\author{Second Author}
\address{Department of Mathematics, University Two, City, Country}

\subjclass[2020]{Primary 05C35; Secondary 11B75}
\keywords{keyword one, keyword two}

\begin{document}

\begin{abstract}
One paragraph: the problem, what you did, what you found, and why it matters. Keep it to about 150 words and avoid citations here.
\end{abstract}

\maketitle

\section{Introduction}

State the question and the main result first.

\begin{theorem}\label{thm:main}
For every integer $n \ge 1$,
\[
  \sum_{k=1}^{n} k = \frac{n(n+1)}{2}.
\]
\end{theorem}

\begin{proof}
By induction on $n$. The case $n = 1$ is immediate, and adding $n+1$ to both
sides of the statement for $n$ gives the statement for $n+1$.
\end{proof}

\begin{remark}
Theorem~\ref{thm:main} is classical; see~\cite{mcmahan2017communication} for a citation example.
\end{remark}

\section{Preliminaries}

\begin{definition}
A \emph{graph} is a pair $G = (V, E)$ where $E \subseteq \binom{V}{2}$.
\end{definition}

\bibliographystyle{amsplain}
\bibliography{references}

\end{document}
